\documentclass[11pt]{article}

% ============================================================
% Physics Exercise Template
% ============================================================

% Mathematics
\usepackage{amsmath}
\usepackage{amssymb}
\usepackage{mathtools}
\usepackage{bm}

% Physics notation
\usepackage{physics}
\usepackage{siunitx}
\usepackage{braket}

% Typography and layout
\usepackage{microtype}
\usepackage[margin=1in]{geometry}
\usepackage{enumitem}

% Figures and references
\usepackage{graphicx}
\usepackage{float}
\usepackage[hidelinks]{hyperref}
\usepackage{cleveref}

% ============================================================
% Basic notation
% ============================================================

\newcommand{\R}{\mathbb{R}}
\newcommand{\C}{\mathbb{C}}
\newcommand{\N}{\mathbb{N}}
\newcommand{\Z}{\mathbb{Z}}

\newcommand{\id}{\mathbb{1}}
\newcommand{\diag}{\operatorname{diag}}
\renewcommand{\tr}{\operatorname{Tr}}

\newcommand{\qed}{\hfill$\square$}

% ============================================================
% Physical symbols
% ============================================================

\newcommand{\lagrdens}{\mathcal{L}}
\newcommand{\nablavec}{\vec{\nabla}}

% ============================================================
% Document information
% ============================================================

\newcommand{\studentname}{Heewon Lee}
\newcommand{\coursecode}{QU.50013}
\newcommand{\coursename}{Quantum Communications and Networks}
\newcommand{\semester}{Fall 2026}


\title{\coursename (\coursecode)\\Exercise Set 2}
\author{\studentname}
\date{\semester}

\begin{document}

\maketitle

\section*{Problem 1}
\noindent
(a)
\[
    \rho_{AB} \overset{\cdot}{=} \begin{pmatrix}
        \cos^2\theta & 0 & 0 & \sin\theta\cos\theta \\
        0 & 0 & 0 & 0 \\
        \sin\theta\cos\theta & 0 & 0 & \sin^2\theta
    \end{pmatrix}
\]
\[
    \Rightarrow \rho_A = \tr_B \rho_{AB} \overset{\cdot}{=} \begin{pmatrix}
        \tr \begin{pmatrix} \cos^2\theta & 0 \\ 0 & 0 \end{pmatrix} &
        \tr \begin{pmatrix} 0 & \sin\theta\cos\theta \\ 0 & 0 \end{pmatrix} \\
        \tr \begin{pmatrix} 0 & 0 \\ \sin\theta\cos\theta & 0 \end{pmatrix} &
        \tr \begin{pmatrix} 0 & 0 \\ 0 & \sin^2\theta \end{pmatrix}
    \end{pmatrix} = \begin{pmatrix}
        \cos^2\theta & 0 \\ 0 & \sin^2\theta
    \end{pmatrix}
\]

\noindent
(b) $\rho_B$ is in its diagonal form in the computational basis,
so its eigenvalues are $\cos^2\theta$ and $\sin^2\theta$.
\[
    \Rightarrow S(\rho_A) = -\cos^2\theta\log_2\cos^2\theta - \sin^2\theta\log_2\sin^2\theta
\]

\noindent
(c) $S(\rho_A)(\theta=0) = 0,\; S(\rho_A)(\theta=\pi/4) = 1$.
In fact, the former corresponds to the classical state $\ket{00}$,
while the latter corresponds to the Bell state
$\ket{\Phi^+} = \frac{1}{\sqrt{2}} (\ket{00} + \ket{11})$.

\noindent
(d) \emph{If the global state is known to be pure}, then the fact that
$\rho_A$ describes a mixed state could allude to global entanglement,
although no local operation nor measurement can distinguish a mixed state
from a subsystem of an entangled state with identical reduced density matrices.

\section*{Problem 2}
\noindent
(a) $E(0) = -1,\; E(\pi/3) = -1/2,\; E(\pi/2) = 0,\; E(\pi) = +1$

\noindent
(b) $\theta=0$ corresponds to perfect anticorrelation, i.e., measurement outcomes are always opposite.
As $\theta$ increases to $\pi/2$, $E$ increases to zero, meaning a net anticorrelation
but with less and less significance.
At $\theta = \pi/2$, $E = 0$ implies the absence of correlation.
After that point, the measurement outcomes are positively correlated,
with more and more significance as $\theta$ increases to $\pi$.
Finally, at $\theta=\pi$, the measurements are perfectly correlated, i.e., always identical outcomes.

\noindent
(c) If Bob measures along $z$, Alice can be absolutely certain that Bob will obtain $-1$.
Along $x$, Alice cannot make any prediction about Bob's outcome.

\noindent
(d) These correlations do not allow superluminal information transfer,
as these correlations can only be compared and utilized once the two measurement outcomes
have been compared at a common point in space.

\end{document}
